\subsection{Self-Trapped Excitons}

Thermal fluctuations induce slight lattice vibrations that act as a trigger, creating a transient local deformation. The strongly localized exciton is captured by this distortion. Through massive exciton-phonon coupling, it abruptly discharges its potential energy into the lattice, considerably deepening the local potential well.

This intense localized energy weakens and breaks the $\text{Si--O--Si}$ covalent bond, displacing the oxygen atom from its equilibrium position. This structural rupture generates two dangling bonds, each holding an unpaired electron. In this strongly distorted configuration, the exciton self-traps. The hole localizes on the oxygen dangling bond, while the electron remains bound to the silicon dangling bond~\cite{mao2004}.

This mechanism explains the formation of \textbf{Self-Trapped Excitons} (STE) in silica, where the excitation energy is rapidly converted into local lattice distortion, opening the way to permanent structural modifications of the material~\cite{mao2004}.
